\documentclass[conference]{IEEEtran}
\usepackage{amsmath}

\begin{document}
\title{Federated Forecasting of Electricity Demand Across Distribution Substations}
\author{\IEEEauthorblockN{Daniel Okafor and Hannah Weiss}\IEEEauthorblockA{Northfield Institute of Technology\\ d.okafor@northfield-tech.edu}}
\maketitle
\begin{abstract}
Distribution operators forecast load per substation, but each utility keeps its meter data private. We train one shared forecaster across 41 substations owned by three utilities using federated averaging, so raw meter readings never leave their owner. The shared model lowers the day-ahead error by 11\% compared with models trained on each utility alone, and costs 3 MB of traffic per training round.
\end{abstract}
\begin{IEEEkeywords}
federated learning, load forecasting, smart grid, privacy
\end{IEEEkeywords}
\section{Introduction}
Day-ahead load forecasts decide how much generation is scheduled and how transformers are loaded. Small utilities have
few substations and short histories, so their forecasters overfit. Pooling data would help, but meter data reveals
when homes are empty, and contracts often forbid sharing it.

Federated learning trains one model across data owners by exchanging model updates instead of data \cite{mcmahan}. We
apply it to substation load, where the owners are few, the series are long, and each owner sees a different climate.

\section{Method}
Each substation forecaster is a two-layer long short-term memory network \cite{hochreiter} that reads the last seven
days of hourly load with temperature and calendar features and predicts the next 24 hours. Every round, each utility
trains the current shared model for one epoch on its own substations with Adam \cite{kingma} and sends back the weight
change. The server averages the changes, weighted by the number of training windows. Dropout \cite{srivastava} on the
recurrent outputs kept the shared model from fitting the largest utility at the expense of the others.

\section{Results}
We used two years of hourly data from 41 substations; the last three months were held out. Table~\ref{t2} shows the
mean absolute percentage error of the day-ahead forecast. The shared model helps the smallest utility the most, whose
seven substations gain 19\%. Training converged after 60 rounds, about 180 MB of traffic per utility in total.

\begin{table}[t]
\caption{Day-ahead error (MAPE)\label{t2}}
\centering
\begin{tabular}{lccc}
\hline
Training & Utility A & Utility B & Utility C \\
\hline
Local only & 6.1\% & 5.4\% & 8.9\% \\
Federated & 5.6\% & 5.1\% & 7.2\% \\
\hline
\end{tabular}
\end{table}

\section{Conclusion}
Utilities can share what their data teaches without sharing the data. Federated averaging closed most of the gap to a
pooled model in our study, and the gain was largest where data was scarcest. Deep sequence models \cite{lecun} are now
cheap enough to run at every substation, which makes this practical for small operators.

\begin{thebibliography}{99}
\bibitem{mcmahan} B. McMahan, E. Moore, D. Ramage, S. Hampson, and B. A. y Arcas, ``Communication-efficient learning of deep networks from decentralized data,'' in \emph{Proc. AISTATS}, 2017, pp. 1273--1282.
\bibitem{hochreiter} S. Hochreiter and J. Schmidhuber, ``Long short-term memory,'' \emph{Neural Computation}, vol. 9, no. 8, pp. 1735--1780, 1997.
\bibitem{kingma} D. P. Kingma and J. Ba, ``Adam: A method for stochastic optimization,'' in \emph{Proc. ICLR}, 2015.
\bibitem{lecun} Y. LeCun, Y. Bengio, and G. Hinton, ``Deep learning,'' \emph{Nature}, vol. 521, pp. 436--444, 2015.
\bibitem{srivastava} N. Srivastava, G. Hinton, A. Krizhevsky, I. Sutskever, and R. Salakhutdinov, ``Dropout: A simple way to prevent neural networks from overfitting,'' \emph{J. Mach. Learn. Res.}, vol. 15, pp. 1929--1958, 2014.
\end{thebibliography}
\end{document}
